% META: {"id":"1SPE-GEOREP-CO-043","chapitre":"1SPE-GEOMETRIE-REPEREE","type_objet":"corrige","exercice_ref":"1SPE-GEOREP-EX-043","capacites_codes":["C5"],"sources_inspiration":[],"status":"approved","origin":"generated","status_history":["generated","audited","accepted","approved"],"release_acceptance":"RELEASE_OWNER_BATCH_ACCEPTANCE","acceptance_closure_digest":"sha256:9b3ccf9a81c5520fbb7e03b7b2d3e7bf2057a02834b6d908bf3be5f49f3b553f","acceptance_receipt_digest":"sha256:4fad86f0282e0fa4e0419cca1ad6379ae7af5a280c04a4a90880f90a1c044242"}
% BEGIN-VERIFY
% from sympy import *
% x, y = symbols('x y')
% # Projete orthogonal de P(4,5) sur D: 3x + 4y - 6 = 0
% # Perpendiculaire: 4x - 3y + c = 0 passant par P: 16-15+c=0 => c=-1
% # Systeme: 3x+4y=6, 4x-3y=1
% sol = solve([3*x+4*y-6, 4*x-3*y-1], [x,y])
% assert sol[x] == Rational(22,25)
% assert sol[y] == Rational(21,25)
% END-VERIFY
\begin{corrige}{1SPE-GEOREP-EX-043}

\commentaireMarge{La perpendiculaire à $\mathcal{D}$ passant par $P$ a pour vecteur normal le vecteur directeur de $\mathcal{D}$.}

Perpendiculaire : $4x - 3y + c = 0$. En $P$ : $16 - 15 + c = 0$, $c = -1$.

$\begin{cases} 3x + 4y = 6 \\ 4x - 3y = 1 \end{cases}$. On multiplie : $9x + 12y = 18$ et $16x - 12y = 4$.

$25x = 22$, $x = \dfrac{22}{25}$. $y = \dfrac{6 - 3 \times 22/25}{4} = \dfrac{84/25}{4} = \dfrac{21}{25}$.

$H\!\left(\dfrac{22}{25} ; \dfrac{21}{25}\right)$.

\end{corrige}
